\chapter{半单模与半单环}

本章中，当我们说到模（不加进一步说明）时，指的是左模。设 A 为环，M 为 A-模。对任意 \(a \in A\) ，我们将 M 的位似 \(x \mapsto ax\) 记作 \(a_{M}\) 。映射 \(a \mapsto a_{M}\) 是从环 A 到一个子环 \(A_{M}\) 的同态，该子环含于 \(\operatorname{End}_{\mathbf{Z}}(\mathrm{M})\) ，我们称之为 M 的位似环。

除非另有说明，我们所考虑的代数都是结合的且含单位元的；所谓代数 E 的子代数，指的是包含 E 的单位元的子代数；代数同态都假设是保单位元的。

设 K 为交换环，L 为交换 K-代数。对任意 K-模 E，我们把由标量扩张得到的 L-模 \(L \otimes_{K} E\) 记作 \(\mathrm{E}_{(\mathrm{L})}\) （II, §5, No.~1, p.~277）。若 A 是 K-代数而 M 是左 A-模，则 \(\mathrm{A}_{(\mathrm{L})}\) 带有一个自然的 L-代数结构（III, §1, No.~5, p.~433），而 \(\mathrm{M}_{(\mathrm{L})}\) 带有左 \(\mathrm{A}_{(\mathrm{L})}\) -模结构，其作用律由公式 \((\lambda \otimes a)(\mu \otimes m) = \lambda \mu \otimes am\) 给出。

对任意交换环 K 和任意群 G，我们把群 G 在环 K 上的代数 \(\mathrm{K}^{(\mathrm{G})}\) 记作 K{[}G{]} （III, §2, No.~6, p.~446）。

